\SourceLeaf{36}

\Needspace{8\baselineskip}
\begin{center}\small
REMARKS ON THE PAPER »BEWEIS DER IRRATIONALITÄT DER TANGENTEN
RATIONALER BÖGEN IN EINER NEUEN GESTALT«.
\end{center}
\par\noindent\hspace*{\parindent}The first sound proof that $\operatorname{tg}\dfrac{m}{n}$ is irrational for integers
$m, n$ was given by \textsc{Legendre} in the fourth note of his »Éléments
de géométrie« (1794). The essential ideas of \textsc{Gauss}'s argument agree with
those of \textsc{Legendre}'s proof. What is new is the way
\textsc{Gauss} proves directly from the power series for the $P_\theta$ that these
quantities continue to decrease as $\theta$ increases. On the investigations
by \textsc{Lambert} and \textsc{Legendre} at issue here, especially on the latter's
account in\,the\,cited work, \textsc{Gauss} wrote a critical
discussion, which is omitted here.

Another development, also omitted here, derives from the
recurrence formulas for the $P_\theta$ approximate representations of $\operatorname{tg} m$
as quotients of polynomial functions of $m$. \textsc{Gauss} uses these
representations to calculate approximations to $\tfrac{1}{2}\pi$.

\begin{flushright}
\textsc{Fricke}.
\end{flushright}
